\clearpage
\chapter*{GENERAL INTRODUCTION}
\addcontentsline{toc}{chapter}{GENERAL INTRODUCTION}

Healthcare has been a fundamental pillar of human civilization, sustaining well-being and driving social stability across the globe. As global populations age, the demand for accessible diagnostics has spurred the adoption of advanced medical technologies, collectively known as HealthTech. Innovations such as Artificial Intelligence (AI), Deep Learning, and data analytics are revolutionizing clinical practices by enabling precision medicine and promoting early disease detection. For instance, global platforms like SkinVision have utilized AI-driven analysis to assess skin cancer risks for millions of users [2], while automated screening systems in Europe have improved diagnostic accuracy by nearly 20% in clinical trials [3]. Digital health applications provide real-time insights into patient symptoms, and platforms like Coursera offer specialized medical training, contributing to a significant rise in healthcare literacy and early intervention rates worldwide [4].

In Tunisia, patients and medical practitioners face significant obstacles that limit the effectiveness of traditional dermatological care. Challenges include a critical shortage of specialized dermatologists, particularly in rural areas, high consultation costs, and long waiting times for clinical appointments [1]. Language barriers further complicate access, as many advanced diagnostic tools are available only in English or French, while a significant portion of the Tunisian population requires localized support. In regions like Sfax, individuals struggle with limited access to preliminary screening tools, resulting in delayed diagnoses for malignant lesions and increased healthcare burdens. These issues underscore the need for a localized, cost-effective, and accessible solution tailored to Tunisia’s healthcare needs, particularly for the early detection of skin cancer.

To address these challenges, the \textbf{DermaAI Platform} was developed in collaboration with \textbf{Aura Dynamic}, a digital innovation agency specializing in AI-driven solutions. This report, titled \textbf{DermaAI Project Report}, documents the project’s development using the \textbf{Agile Scrum} methodology. \textbf{Chapter 1} provides a comprehensive overview, covering the collaboration, global healthcare trends, project objectives, and methodology. \textbf{Chapter 2} details the technical design, including deep learning algorithms, system architecture, and the software environment. \textbf{Chapter 3} describes the implementation of user authentication and access management. \textbf{Chapter 4} presents the core AI integration and skin lesion analysis. \textbf{Chapter 5} covers the scan history, reporting, and chatbot features, while \textbf{Chapter 6} demonstrates the administrative dashboard and system analytics. Aimed at researchers, medical practitioners, and technology developers, this work aims to advance digital dermatology and inspire accessible healthcare practices in Tunisia and beyond.
\newpage
